\section{Background}

\subsection{Programming Languages as Cultural Artifacts}

Programming languages are often described through formal properties: syntax,
semantics, type systems, execution models, module systems, and implementation
techniques. Those properties matter, but they do not exhaust what languages do
in practice. A programming language also supplies a vocabulary for teaching,
debugging, collaboration, and self-identification. Keywords, library names,
error messages, tutorials, and examples shape what a language feels like and
who can imagine themselves as a participant in its community.

% TODO: Cite work in software studies, human-computer interaction, and
% sociotechnical systems that treats programming languages and developer tools
% as cultural artifacts rather than only formal calculi.

\subsection{English Dominance in Programming}

English occupies a dominant position in programming-language syntax, software
documentation, package ecosystems, and error reporting. This dominance is
partly historical and partly infrastructural: early computing systems,
keyboard layouts, standards, and major commercial and academic institutions
privileged English-language naming. For programmers outside English-dominant
contexts, this creates a layered literacy problem. A learner must acquire both
computational concepts and a specialized English-derived vocabulary for those
concepts.

% TODO: Add citations on English dominance in computing, global software
% development, and language barriers in programming education.

\subsection{Educational Languages, Esolangs, and Localization}

Localized programming environments have appeared in several forms. Some
educational systems translate commands, blocks, or messages so that beginners
can focus on concepts before encountering English-heavy professional tooling.
Other languages use non-English syntax as an artistic, political, humorous, or
community-building gesture. Esoteric languages also show that programming
language design can carry cultural expression, but they are often evaluated by
different standards than educational or practical languages.

\plat{} sits between these categories. It is small enough to resemble a
beginner-oriented teaching language, expressive enough to run scripts with
variables, functions, tables, loops, and conditionals, and linguistically
marked enough to invite comparison with esolangs and cultural programming
artifacts. The project therefore provides an opportunity to ask how a localized
language can be playful or expressive without being reduced to novelty.

% TODO: Replace the placeholder citations in references.bib with accurate
% sources on localized programming languages, block-based educational tools,
% esolangs, and non-English programming environments.
